\section{Sissejuhatus} \label{sissejuhatus}

Suured keelemudelid (\textit{Large Language Models}, LLM) on kahe-kolme aasta jooksul liikunud tehnoloogiaentusiastide mänguasjast igapäevaseks töövahendiks, mida kasutatakse tõsisemate tekstide koostamisel nii ettevõtetes kui ka kõrgkoolides \cite{liang2024mapping}. OpenAI GPT-seeria, Google'i Gemini ning Anthropicu Claude'i mudelid suudavad tänaseks koostada sidusat akadeemilist teksti, toimetada olemasolevaid dokumente ja genereerida ka pikemaid struktureeritud kirjalikke töid \cite{brown2020language, openai2024gpt4}. Väiksemate keelte -- sealhulgas eesti keele -- puhul on mudelite kvaliteet aga ebaühtlasem ning selle süstemaatiline hindamine on seni jäänud tagasihoidlikuks \cite{kuulmets2025finno, dorkin2026estllm}.

Bakalaureusetöö on Tartu Ülikooli arvutiteaduse instituudis (ATI) esimene suurem iseseisev akadeemiline kirjatöö, mille üliõpilane peab vormistama kehtestatud juhendi kohaselt \cite{tuati2025juhend}. Töö valmimine nõuab tavaliselt ühe semestri, mille käigus tuleb süveneda kirjandusse, kirjeldada metoodikat, esitada tulemused ja järgida akadeemilist stiili. Üliõpilaste seas on viimastel aastatel laialdaselt levinud TI mudelite kasutamine lõputöö kirjutamise abivahendina, kuid senised juhendid ja uurimused ei anna selget pilti sellest, kui usaldusväärselt TI mudelid sellise ülesandega eesti keeles toime tulevad \cite{lund2025perceptions, rodafinos2025integration}. See küsimus on nii praktiliselt oluline (üliõpilane vajab tõenduspõhist otsust tööriista valikul) kui ka akadeemiliselt oluline (juhendajad ja ülikoolid peavad kujundama selgeid reegleid TI kasutamiseks).

Käesoleva bakalaureusetöö eesmärk on eksploratiivselt hinnata, kas ja kuidas tuleksid kaasaegsed TI mudelid koos tänapäevaste agentsete tööriistadega toime eestikeelse tehnilise informaatika lõputöö genereerimisega. Töö keskendub järgmistele uurimisküsimustele:

\begin{enumerate}
    \item Millised kaasaegsete TI mudelite ja agentsete tööriistade kombinatsioonid tulevad toime mahukate eestikeelsete akadeemiliste tekstide genereerimisega?
    \item Kuidas erinevad nende kombinatsioonide väljundid, kui kõigile antakse üks ja sama lühike kasutajaviibe -- kas ja kuidas erinevad need mudel-tööriista paarid ülesande, keelevaliku ja iseseisvalt valitud teema tajumisel?
    \item Millistes lõputöö osades suudab TI pakkuda tõhusat abi ning kus on inimese panus endiselt asendamatu?
    \item Millised on TI-põhise tekstigenereerimise praktilised piirangud, sealhulgas hinnapoliitika, kontekstiakna piirangud ja keelelised puudujäägid?
\end{enumerate}

Uurimuse läbiviimiseks valiti neli juhtivat suurt keelemudelit -- Google Gemini 3.1 Pro, OpenAI GPT-5.4, Anthropic Claude Sonnet 4.6 ning Anthropic Claude Opus 4.6 -- ja neli nende kasutamise viisi: Gemini 3.1 Pro koos Google'i Antigravity agentse keskkonnaga, GPT-5.4 nii tavapärase ChatGPT vestlusliidese kui ka ChatGPT Codex laienduse (Visual Studio Code'is) kaudu ning Claude Sonnet 4.6 ja Opus 4.6 koos Anthropicu Cowork agentse tööriistaga (Cowork on sama ettevõtte toodetud Claude Code'ist eraldiseisev agentne tööriist). Igale mudel-tööriista kombinatsioonile anti üks ja sama lühike eestikeelne kasutajaviibe, mis palus mudelil genereerida LaTeX-vormingus terviklik bakalaureusetöö Tartu Ülikooli arvutiteaduse instituudi nõuetele vastavalt, valida ise informaatikaga seotud teema ning kasutada töökausta paigutatud UniTartuCS LaTeX-malli. Iteratiivseid järelküsimusi, peatükkide kaupa juhendamist ega juhendite või näidistööde eraldi sisestamist konteksti ei kasutatud; kogu mudelitele antud juhis mahtus ühte lühikesse viipesse. Genereeritud tervikdokumentide kvaliteeti hinnati mitme kriteeriumi alusel, sealhulgas struktuuri, kirjavigade tiheduse, viidete hulga ja kontrollitavuse, eestikeelse akadeemilise stiili (anglitsismid, esimese isiku asesõnade kasutus) ning mahu põhjal. Kuna iga kombinatsiooni kohta tehti vaid üks jooks, on saadud tulemusi käsitletud eksploratiivselt; suurte keelemudelite mitteterministlikkuse tõttu võivad korduskatsed anda erinevaid tulemusi ning järeldused ei pretendeeri üldistatavusele kogu mudeli võimekusele (vt peatükk~\ref{subsec:puudujaagid}).

Töö on struktureeritud järgmiselt. Peatükk~\ref{taust} annab ülevaate suurte keelemudelite arengust, eesti keele töötlemise väljakutsetest ning varasematest uurimustest TI kasutamisest akadeemilises kirjutamises. Peatükk~\ref{metoodika} kirjeldab uurimuse metoodikat, sealhulgas mudelite valiku kriteeriume, ühe kasutajaviibega agentse eksperimendi ülesehitust ja hindamiskriteeriumeid. Peatükk~\ref{mudelid} esitab valitud mudelite detailse võrdluse, hõlmates nende tehnilist arhitektuuri, hinnastamist ja eesti keele tuge. Peatükk~\ref{eksperimendid} kirjeldab läbiviidud eksperimenti ja selle tulemusi -- tervikliku bakalaureusetöö genereerimise katset ühe kasutajaviibega, mis paljastas mudel-tööriista kombinatsioonide vahel põhimõttelisi erinevusi juhiste järgimisel, teemavalikul ja keelevalikul. Peatükk~\ref{arutelu} sisaldab arutelu tulemuste üle, sealhulgas TI kasutamise praktilised soovitused, piirangud ja tulevikuväljavaated. Peatükk~\ref{claudecode} kirjeldab Claude Code'i agentset tööriista ja sellega sarnast Anthropicu toodet Cowork ning esitab praktilised soovitused lõputöö genereerimiseks. Peatükk~\ref{kokkuvote} võtab töö tulemused kokku.
